\subsection{TOPOLOGY OF REAL PROJECTIVE SPACES}

We recall the following definitions from algebra: given a division ring or a field \(\mathbf{K}\) , the set of lines passing through \(o\) (i.e.~vector subspaces of dimension 1) in the left vector space \(\mathbf{K}_s^{n+1}\) over \(\mathbf{K}\) is called left projective space of \(n\) dimensions over \(\mathbf{K}\) , and is denoted by \(\mathbf{P}_n(\mathbf{K})\) .

If we make correspond to every line passing through o in \(K_{s}^{n+1}\) the same line with the origin omitted, we have a bijection of \(\mathbf{P}_{n}(\mathbf{K})\) onto the quotient of \(K_{n+1}^{*}\) (the complement of \(\{o\}\) in \(K_{s}^{n+1}\) ) by the following equivalence relation \(\Delta_{n}(\mathbf{K})\) between vectors x and y of \(K_{n+1}^{*}\) : ``there exists \(t \in K\) such that \(t \neq o\) and y = tx''. In what follows we shall identify \(P_{n}(\mathbf{K})\) with this quotient set. In the theory of projective spaces, we take the interval \([o, n]\) of N as index set for the coordinates of a point of \(K_{n+1}^{*}\) . The coordinates \(x_{i}\) ( \(o \leqslant i \leqslant n\) ) of any one of the points of \(K_{n+1}^{*}\) whose canonical image is \(x \in P_{n}(\mathbf{K})\) constitute what is called a system of homogeneous coordinates of the point x.

For each integer \(p\) such that \(-1 \leqslant p \leqslant n\) , the canonical image in \(\mathbf{P}_n(\mathbf{K})\) of a vector subspace of \(p + 1\) dimensions (without the origin) of \(\mathrm{K}_s^{n+1}\) is called a projective linear variety of \(p\) dimensions. A system of \(p\) points of \(\mathbf{P}_n(\mathbf{K})\) is said to be free if it consists of the canonical images of \(p\) points of \(\mathrm{K}_{n+1}^*\) which form a free system in the vector space \(\mathrm{K}_s^{n+1}\) . The projective linear variety of \(\mathbf{P}_n(\mathbf{K})\) generated by a free system of \(p + 1\) points (i.e.~the smallest projective linear variety which contains these \(p + 1\) points) has dimension \(p\) .

When K is the field R, the corresponding projective spaces can be topologized, and it is these we shall study.

DEFINITION 1. The projective space \(\mathbf{P}_n(\mathbf{R})\) endowed with the quotient topology of the topology of \(\mathbf{R}_{n+1}^*\) by the equivalence relation \(\Delta_n(\mathbf{R})\) is called real projective space of \(n\) dimensions.

The projective space \(\mathbf{P}_{1}(\mathbf{R})\) is called the real projective line, and \(\mathbf{P}_{2}(\mathbf{R})\) is called the real projective plane.

Whenever there is no risk of confusion, we shall write \(P_{n}\) and \(\Delta_{n}\) instead of \(\mathbf{P}_{n}(\mathbf{R})\) and \(\Delta_{n}(\mathbf{R})\) .

PROPOSITION 1. The projective space \(P_{n}\) is Hausdorff.

We start by showing that the relation \(\Delta_{n}\) is open (Chapter I, § 5, no. 2). Let A be an open set in \(\mathbf{R}_{n+1}^{*}\) ; to saturate A with respect to \(\Delta_{n}\) we have to take the union of the sets \(t\mathbf{A}\) homothetic to A, as \(t\) runs through the set of real numbers \(\neq 0\) ; since each of these sets is open, so is their union.

By Proposition 8 of Chapter I, § 8, no. 3, the proposition will be proved if we show that the subset M of \(\mathbf{R}_{n+1}^{*}\times\mathbf{R}_{n+1}^{*}\) defined by the relation \(\Delta_{n}\) is closed. Let then \((x,y)\) be a point of \(\mathbf{R}_{n+1}^{*}\times\mathbf{R}_{n+1}^{*}\) lying in the closure of M. If \(x=(x_{i})\) , there is an index \(i\) such that \(x_{i}\neq0\) ; hence there is a neighbourhood V of \((x,y)\) such that for every point \((x^{\prime},y^{\prime})\in\mathbf{M}\cap\mathbf{V}\) the \(i\) th coordinate \(x_{i}^{\prime}\) of \(x^{\prime}\) is not o. As \((x^{\prime},y^{\prime})\) tends to \((x,y)\) while remaining in M, \(y_{i}^{\prime}x_{i}^{\prime-1}\) tends to \(t=y_{i}x_{i}^{-1}\) ; since \(y^{\prime}=(y_{i}^{\prime}x_{i}^{\prime-1})x^{\prime}\) , we see by passing to the limit that \(y=tx\) , which shows that \((x,y)\in\mathbf{M}\) .

PROPOSITION 2. The projective space \(\mathbf{P}_{n}\) is compact and connected, and is homeomorphic to the quotient of the sphere \(\mathbf{S}_{n}\) by the equivalence relation induced on the sphere by \(\Delta_{n}\) .

Let \(\Delta_n'\) be the equivalence relation induced on \(\mathbf{S}_n\) by \(\Delta_n\) (the equivalence classes of \(\Delta_n'\) are pairs of diametrically opposite points of \(\mathbf{S}_n\) ). The mapping \(x \to x/||x||\) of \(\mathbf{R}_{n+1}^*\) onto \(\mathbf{S}_n\) is continuous, hence (Lemma 2) \(\mathbf{P}_n\) is homeomorphic to \(\mathbf{S}_n/\Delta_n'\) . Since \(\mathbf{S}_n\) is compact and connected, every Hausdorff quotient space of \(\mathbf{S}_n\) is also compact and connected (Chapter I, § 9, no. 4, Theorem 2, Corollary 1; § 11, no. 3, Proposition 7).

PROPOSITION 3. If \(n \geqslant 0\) , the projective space \(\mathbf{P}_n\) is homeomorphic to the quotient space of the ball \(\mathbf{B}_n\) obtained by identifying each point of \(\mathbf{S}_{n-1}\) with its diametrically opposite point.

Let H be the closed hemisphere of \(S_{n}\) defined by \(x_{0} \leqslant 0\) . \(P_{n}\) , which is homeomorphic to the quotient space of \(S_{n}\) by the relation \(\Delta_{n}^{\prime}\) , is also homeomorphic to the quotient of the subset H of \(S_{n}\) by the relation \(\Delta_{n}^{\prime\prime}\) induced on H by \(\Delta_{n}^{\prime}\) . For every equivalence class of \(\Delta_{n}^{\prime}\) meets H in at least one point, and therefore (Lemma 1) it is enough to verify that if we saturate with respect to \(\Delta_{n}^{\prime}\) an open subset U of H which is saturated with respect to \(\Delta_{n}^{\prime\prime}\) , we get an open subset V of \(S_{n}\) . Now if \(a = (a_{i}) \in U\) and if \(a_{0} < 0\) , there is a neighbourhood W of a in \(S_{n}\) contained in U, and the union of W and \(-W\) is a neighbourhood of a saturated with respect to \(\Delta_{n}^{\prime}\) and contained in V. If on the other hand \(a_{0} = 0\) , we have \(- a \in U\) , and there exists r \textgreater{} o such that the set of points \(x \in H\) which satisfy one or the other of the relations \(||x-a||<r, ||x+a||<r\) is contained in U; the set of points \(x\in S_{n}\) which satisfy one or the other of these relations is the neighbourhood of a which is saturated with respect to \(\Delta_{n}^{\prime}\) and is contained in V.

Observe that the quotient space \(H/\Delta_{n}^{\prime\prime}\) is obtained by identifying, in H, each point of the intersection \(S_{n-1}\) of H and the hyperplane \(x_{0}=0\) with its opposite point. To complete the proof it suffices to remark that the stereographic projection with vertex \(e_{0}\) (§ 2, no. 4) is a homeomorphism of H onto \(B_{n}\) which leaves invariant the points of \(S_{n-1}\) .

